% Transcrição do PDF de Natanael Henrik Zago, páginas impressas 23--25.
\chapter{TRABALHOS RELACIONADOS}
\label{chap:relacionados}
\section{Teste de desvio: uma geração de casos de teste - Técnica para APIs GraphQL}

Vargas et al. \cite{vargas2018} criaram uma ferramenta para testar APIs GraphQL de forma automática explorando a abordagem de desvio, ou seja, a partir de uma consulta GraphQL informada que fosse válida, a ferramenta gera novos casos de testes, variações da consulta original. Após gerados esses novos casos, eram realizados os testes e verificado se os resultados obtidos das novas variações de casos de teste eram um subconjunto do caso de teste original, ou mesmo se continha erros.

Além de desenvolver a ferramenta Vargas et al. \cite{vargas2018}, foram realizados 3 estudos de caso para ver a capacidade da ferramenta em encontrar falhas. Ao analisar o resultado dos estudos de caso foi constatado que a ferramenta foi capaz de encontrar casos de erro no retorno das consultas.

Como principais diferenças em relação ao trabalho de Vargas et al. \cite{vargas2018}, o presente trabalho vai realizar a geração das consultas baseado em especificação, enquanto o trabalho de \cite{vargas2018} faz a geração das consultas através de desvios de uma consulta original pré-estabelecida.

\section{AutoGraphQL: Uma ferramenta de geração de teste automatizada para GraphQL}
Zetterlund \cite{zetterlund2021} implementou uma ferramenta para realização de teste de APIs GraphQL chamada AutoGraphQL, esta ferramenta funciona da seguinte forma: ela coleta consultas realizadas na API em produção e armazena em um banco de dados, para posteriormente utilizar esses dados como entradas para as consultas geradas baseadas no esquema da API. Como tecnologias usadas para implementação do AutoGraphQL, foi utilizado a linguagem Python e os casos de teste gerados para validação das consultas com o PHPunit, este último utilizado para validar se os dados retornados nas consultas possuíam a tipagem correta.

No trabalho de Zetterlund \cite{zetterlund2021}, a criação da ferramenta AutoGraphQL foi utilizada para um estudo de caso em específico, onde foram coletados dados de consultas durante 33 dias, e posteriormente efetuado os testes na API GraphQL em análise. Como resultados obtidos, eles conseguiram encontrar 157 falhas na API num total de 24.049 de casos de testes gerados.

O trabalho proposto neste projeto de pesquisa se diferencia do trabalho de Zetterlund \cite{zetterlund2021} por fazer a geração de consultas de forma aleatória baseado no esquema da API, ou seja, não necessita de consultas previas para criação dos casos de teste, bem como as tecnologias utilizadas para elaboração da ferramenta que neste trabalho será utilizado a linguagem funcional Haskell.

\section{Teste automático baseado em propriedade de APIs GraphQL}
Karlsson; Causevic; Sundmark \cite{karlsson2020graphql} apresentam um dos primeiros trabalhos que propõe a geração de consultas GraphQL de forma aleatória para testes de API GraphQL. Uma vez criada as consultas são realizados testes na API e então verificado o retorno para apuração se a resposta da API está em conformidade com a definição de tipo especificados no esquema. No mesmo trabalho foi proposto observar a cobertura de casos de teste atingido pelas consultas geradas de forma aleatória, bem como analisar a aplicação da ferramenta desenvolvida para um estudo de caso.

Como resultado do trabalho de Karlsson; Causevic; Sundmark \cite{karlsson2020graphql} foi observado que foi possível atingir uma grande cobertura de casos de teste, em algumas configurações até em 100\%. No estudo de caso também foi possível observar que a ferramenta foi capaz de encontrar falhas na API em análise de teste, atendendo assim com seu propósito.

As principais distinções propostas no presente trabalho está na tecnologia utilizada para geração das consultas utilizando a linguagem funcional Haskell diferente da linguagem Clojure usada no trabalho de Karlsson; Causevic; Sundmark \cite{karlsson2020graphql}, bem como o foco do trabalho que é criar um gerador de consultas e utilizar o teste baseado em propriedades para analisar as consultas geradas, já no trabalho de Karlsson; Causevic; Sundmark \cite{karlsson2020graphql} não foi adentrado em como são geradas as consultas, mas sim a averiguação da propriedade de tipos depois de já executado uma consulta na API.

\section{Fuzzing de caixa branca e caixa preta para APIs GraphQL}
Belhadi; Zhang; Arcuri \cite{belhadi2022} desenvolveram uma ferramenta para teste de APIs GraphQL que usa duas abordagens, testes de caixa branca, quando se possui acesso ao código fonte e testes de caixa preta quando não se possui acesso ao código fonte, onde os casos de teste de caixa branca utilizam uma abordagem evolutiva e para os testes de caixa preta uma abordagem aleatória.

No trabalho foi ressaltado uma maior precisão ao se utilizar teste de caixa branca, porém os testes de caixa preta apesar de apresentarem menor exatidão, também foram capazes de encontrar falhas nas APIs em teste, bem como tem vantagem de serem aplicados quando não se possui o acesso ao código fonte.

O trabalho de Belhadi; Zhang; Arcuri \cite{belhadi2022} não teve enfoque na geração de consultas GraphQL, mas sim na execução de casos de testes na API para verificar a resposta e analisar possíveis inconformidades, isto para os casos de teste de caixa preta. Já para o casos de testes de caixa branca foi realizado uma análise sobre o próprio código fonte da API, que o trabalho aqui proposto não pretende adentrar.

\section{QuickREST: Geração de teste baseada em propriedade de APIs RESTful descritas pela OpenAPI}
Stefan \cite{karlsson2020rest} apresenta uma ferramenta para automatizar testes de APIs RESTful, utilizando como base para os testes a documentação da API descrita pela OpenAPI, que a partir da descrição é realizada a geração dos casos de teste baseados em propriedades de tipos. O trabalho visa diminuir o trabalho humano para a criação dos testes e realização dos mesmos, bem como aumentar a confiabilidade das APIs RESTful onde a ferramenta desenvolvida foi aplicada.

Além de desenvolver a ferramenta QuickREST, Stefan \cite{karlsson2020rest} realizou um estudo de caso da aplicação da ferramenta em um caso da indústria. Com o estudo de caso foi possível constatar que a ferramenta foi capaz de encontrar erros na API testada, assim mostrando sua utilidade para testes de API.

O trabalho de Stefan \cite{karlsson2020rest} funciona exclusivamente para APIs RESTful, não permitindo testes para outros modelos de API, como por exemplo APIs GraphQL, que é um modelo de exposição de APIs que vem se difundindo nos últimos tempos, e será proposto no presente projeto de pesquisa.
